% Cycle itératif et incrémental appliqué au projet (figure 1.3)
\begin{tikzpicture}[x=1cm, y=1cm]

  \def\R{2.55}

  % Six étapes réparties sur le cercle, en partant du haut
  \foreach \i/\lbl in {
      0/{Analyse\\du besoin},
      1/{Conception},
      2/{Réalisation},
      3/{Mesure},
      4/{Conservation\\ou reprise},
      5/{Intégration de\\l'incrément}}
    {
      \pgfmathsetmacro{\ang}{90 - \i*60}
      \node[fbox, minimum width=21mm, minimum height=9mm]
        (e\i) at ({\ang}:\R) {\lbl};
    }

  % Flèches circulaires entre étapes consécutives
  \foreach \i in {0,1,2,3,4}
    {
      \pgfmathtruncatemacro{\j}{\i+1}
      \draw[flux] (e\i) to[bend right=18] (e\j);
    }
  \draw[flux] (e5) to[bend right=18] (e0);

  % Noyau : l'incrément
  \node[fboxgg, circle, minimum size=15mm, align=center, font=\sffamily\scriptsize\bfseries]
    (noyau) at (0,0) {Incrément\\numéroté};

  % Ce que porte chaque incrément, en note latérale
  \node[note, text width=34mm, anchor=west] (n1) at (\R+1.6, 0.42)
    {Une branche dédiée\\Un document de version\\Un changement réversible};

  \node[note, text width=34mm, anchor=west] (n2) at (\R+1.6, -1.35)
    {La mesure décide :\\l'incrément est conservé,\\corrigé ou remplacé};

  \draw[dependance] (noyau.east) -- (n1.west);
  \draw[dependance] (e3.east) to[bend right=8] (n2.west);

\end{tikzpicture}
